% GENERATED FILE. Edit canonical lessons, then run npm run generate:books.
% canonical-source-hash: fnv1a64:e89ac710381db5d1
% canonical-lessons: UR-C490-taya, UR-C490-khalu, UR-C490-sautela, UR-C490-yatim, UR-C490-god-lena

\chapter{Uncles, Steps and Orphans}
\label{ch:ur-uncles-steps-and-orphans}
\hlchaptermodality{490}

\begin{chapteropening}
\textbf{By the end of this chapter:} \emph{I can talk about uncles, step-parents and orphans.}

Complete the last lesson of chapter 490: I can talk about uncles, step-parents and orphans.
\end{chapteropening}

\section[\texorpdfstring{\ur{تایا} (tāyā)}{تایا (tāyā)}]{\ur{تایا} (tāyā) — a father's elder brother}

\label{lesson:UR-C490-taya}



\begin{quote}
Before the new one: say the Urdu for a nephew (a sister's son), then the Urdu for a grandson (a daughter's son).
\end{quote}

\subsection*{You'll want to know: \ur{تایا}}
\textbf{\ur{تایا}} — \emph{tāyā} — \textquotedblleft{}a father's elder brother\textquotedblright{}.

\begin{quote}
\ur{میرے} \ur{تایا} \ur{ابو} \ur{سے} \ur{بڑے} \ur{ہیں۔}
\end{quote}

\begin{quote}
\emph{mere tāyā abbū se baṛe haiṅ.}
\end{quote}

\begin{quote}
My uncle is older than my father.
\end{quote}

Masculine.

\begin{cousinweb}
\textbf{\ur{تایا}} (\emph{tāyā}) is inherited: CDIAL lists the Hindi cousin \emph{tāyā} under Sanskrit \emph{tātyá}, \textquotedblleft{}fatherly\textquotedblright{} (CDIAL 5759).
\end{cousinweb}

\begin{grammarlens}[title={What you've built}]
One more word for describing what you lived and felt.
\end{grammarlens}

\subsection*{Your turn}
\begin{itemize}
\raggedright
  \item \emph{Say it:} \emph{tāyā}
  \item \emph{Say it:} \emph{tāyā}, once more
  \item \emph{Say it:} \emph{mere tāyā abbū se baṛe haiṅ.} — My uncle is older than my father.
  \item \emph{Recall:} say the Urdu for a nephew (a sister's son), then the Urdu for a grandson (a daughter's son), then say \emph{tāyā} again
\end{itemize}

\begin{tcolorbox}[breakable,colback=teal!4,colframe=teal!35!black,title={Before you move on}]
What does \ur{تایا} mean? (\textquotedblleft{}A father's elder brother\textquotedblright{}.) Say it once more.
\end{tcolorbox}
\section[\texorpdfstring{\ur{خالو} (khālū)}{خالو (khālū)}]{\ur{خالو} (khālū) — a mother's sister's husband}

\label{lesson:UR-C490-khalu}



\begin{quote}
Before the new one: say the Urdu for a grandson (a daughter's son), then the Urdu for a father's elder brother.
\end{quote}

\subsection*{You'll want to know: \ur{خالو}}
\textbf{\ur{خالو}} — \emph{khālū} — \textquotedblleft{}a mother's sister's husband\textquotedblright{}.

\begin{quote}
\ur{خالہ} \ur{اور} \ur{خالو} \ur{کل} \ur{آئیں} \ur{گے۔}
\end{quote}

\begin{quote}
\emph{khālā aur khālū kal āeṅ ge.}
\end{quote}

\begin{quote}
My aunt and uncle will come tomorrow.
\end{quote}

Masculine.

\begin{cousinweb}
\textbf{\ur{خالو}} (\emph{khālū}) is built on \textbf{\ur{خالہ}}, \textquotedblleft{}a mother's sister\textquotedblright{}, usually traced, through Persian, to Arabic \emph{khāla}: the aunt's husband.
\end{cousinweb}

\begin{grammarlens}[title={What you've built}]
One more word for describing what you lived and felt.
\end{grammarlens}

\subsection*{Your turn}
\begin{itemize}
\raggedright
  \item \emph{Say it:} \emph{khālū}
  \item \emph{Say it:} \emph{khālū}, once more
  \item \emph{Say it:} \emph{khālā aur khālū kal āeṅ ge.} — My aunt and uncle will come tomorrow.
  \item \emph{Recall:} say the Urdu for a grandson (a daughter's son), then the Urdu for a father's elder brother, then say \emph{khālū} again
\end{itemize}

\begin{tcolorbox}[breakable,colback=teal!4,colframe=teal!35!black,title={Before you move on}]
What does \ur{خالو} mean? (\textquotedblleft{}A mother's sister's husband\textquotedblright{}.) Say it once more.
\end{tcolorbox}
\section[\texorpdfstring{\ur{سوتیلا} (sautelā)}{سوتیلا (sautelā)}]{\ur{سوتیلا} (sautelā) — step- (related by remarriage)}

\label{lesson:UR-C490-sautela}



\begin{quote}
Before the new one: say the Urdu for a father's elder brother, then the Urdu for a mother's sister's husband.
\end{quote}

\subsection*{You'll want to know: \ur{سوتیلا}}
\textbf{\ur{سوتیلا}} — \emph{sautelā} — \textquotedblleft{}step- (related by remarriage)\textquotedblright{}.

\begin{quote}
\ur{اس} \ur{کی} \ur{سوتیلی} \ur{ماں} \ur{بڑی} \ur{مہربان} \ur{ہے۔}
\end{quote}

\begin{quote}
\emph{us kī sautelī māṅ baṛī mehrbān hai.}
\end{quote}

\begin{quote}
His stepmother is very kind.
\end{quote}

\begin{cousinweb}
\textbf{\ur{سوتیلا}} (\emph{sautelā}) is inherited: it is built on \emph{saut}, \textquotedblleft{}a co-wife\textquotedblright{}, which CDIAL lists under Sanskrit \emph{sapátnī}, \textquotedblleft{}co-wife\textquotedblright{} (CDIAL 13130).
\end{cousinweb}

\begin{grammarlens}[title={What you've built}]
One more word for describing what you lived and felt.
\end{grammarlens}

\subsection*{Your turn}
\begin{itemize}
\raggedright
  \item \emph{Say it:} \emph{sautelā}
  \item \emph{Say it:} \emph{sautelā}, once more
  \item \emph{Say it:} \emph{us kī sautelī māṅ baṛī mehrbān hai.} — His stepmother is very kind.
  \item \emph{Recall:} say the Urdu for a father's elder brother, then the Urdu for a mother's sister's husband, then say \emph{sautelā} again
\end{itemize}

\begin{tcolorbox}[breakable,colback=teal!4,colframe=teal!35!black,title={Before you move on}]
What does \ur{سوتیلا} mean? (\textquotedblleft{}Step- (related by remarriage)\textquotedblright{}.) Say it once more.
\end{tcolorbox}
\section[\texorpdfstring{\ur{یتیم} (yatīm)}{یتیم (yatīm)}]{\ur{یتیم} (yatīm) — an orphan}

\label{lesson:UR-C490-yatim}



\begin{quote}
Before the new one: say the Urdu for a mother's sister's husband, then the Urdu for related by remarriage.
\end{quote}

\subsection*{You'll want to know: \ur{یتیم}}
\textbf{\ur{یتیم}} — \emph{yatīm} — \textquotedblleft{}an orphan\textquotedblright{}.

\begin{quote}
\ur{وہ} \ur{یتیم} \ur{بچہ} \ur{چچا} \ur{کے} \ur{گھر} \ur{رہتا} \ur{ہے۔}
\end{quote}

\begin{quote}
\emph{voh yatīm bachchā chachā ke ghar rahtā hai.}
\end{quote}

\begin{quote}
That orphan boy lives at his uncle's house.
\end{quote}

Masculine or feminine, by the person meant.

\begin{cousinweb}
\textbf{\ur{یتیم}} (\emph{yatīm}) is usually traced, through Persian, to Arabic \emph{yatīm}, \textquotedblleft{}an orphan\textquotedblright{}.
\end{cousinweb}

\begin{grammarlens}[title={What you've built}]
One more word for describing what you lived and felt.
\end{grammarlens}

\subsection*{Your turn}
\begin{itemize}
\raggedright
  \item \emph{Say it:} \emph{yatīm}
  \item \emph{Say it:} \emph{yatīm}, once more
  \item \emph{Say it:} \emph{voh yatīm bachchā chachā ke ghar rahtā hai.} — That orphan boy lives at his uncle's house.
  \item \emph{Recall:} say the Urdu for a mother's sister's husband, then the Urdu for related by remarriage, then say \emph{yatīm} again
\end{itemize}

\begin{tcolorbox}[breakable,colback=teal!4,colframe=teal!35!black,title={Before you move on}]
What does \ur{یتیم} mean? (\textquotedblleft{}An orphan\textquotedblright{}.) Say it once more.
\end{tcolorbox}
\section[\texorpdfstring{\ur{گود} \ur{لینا} (god lenā)}{گود لینا (god lenā)}]{\ur{گود} \ur{لینا} (god lenā) — to adopt (a child)}

\label{lesson:UR-C490-god-lena}



\begin{quote}
Before the new one: say the Urdu for related by remarriage, then the Urdu for an orphan.
\end{quote}

\subsection*{You'll want to know: \ur{گود} \ur{لینا}}
\textbf{\ur{گود} \ur{لینا}} — \emph{god lenā} — \textquotedblleft{}to adopt (a child)\textquotedblright{}.

\begin{quote}
\ur{انہوں} \ur{نے} \ur{ایک} \ur{یتیم} \ur{لڑکی} \ur{کو} \ur{گود} \ur{لیا۔}
\end{quote}

\begin{quote}
\emph{unhoṅ ne ek yatīm laṛkī ko god liyā.}
\end{quote}

\begin{quote}
They adopted an orphan girl.
\end{quote}

Past, with \ur{نے} (he or she …): \ur{اس} \ur{نے} \ur{بچے} \ur{کو} \ur{گود} \ur{لیا}, \emph{us ne bachche ko god liyā}.

\begin{cousinweb}
\textbf{\ur{گود} \ur{لینا}} (\emph{god lenā}) is built from \textbf{\ur{گود}}, \textquotedblleft{}a lap\textquotedblright{}, inherited from \emph{gōddī}, a form CDIAL reconstructs (CDIAL 4284), and \textbf{\ur{لینا}}, \textquotedblleft{}to take\textquotedblright{}: to take a child into one's lap.
\end{cousinweb}

\begin{grammarlens}[title={What you've built}]
One more word for describing what you lived and felt.
\end{grammarlens}

\subsection*{Your turn}
\begin{itemize}
\raggedright
  \item \emph{Say it:} \emph{god lenā}
  \item \emph{Say it:} \emph{god lenā}, once more
  \item \emph{Say it:} \emph{unhoṅ ne ek yatīm laṛkī ko god liyā.} — They adopted an orphan girl.
  \item \emph{Recall:} say the Urdu for related by remarriage, then the Urdu for an orphan, then say \emph{god lenā} again
\end{itemize}

\begin{tcolorbox}[breakable,colback=teal!4,colframe=teal!35!black,title={Before you move on}]
What does \ur{گود} \ur{لینا} mean? (\textquotedblleft{}To adopt (a child)\textquotedblright{}.) Say it once more.
\end{tcolorbox}
